% Adición a las conclusiones de VII; las conclusiones anteriores se conservan.
\paragraph{Intégration des quatre interactions et clôture canonique totale.}
La gravitation et les interactions forte, faible et électromagnétique agissent sur une même réalisation matérielle de la structure APP--TRIT--TPK. Le transport conserve les feuillets, l’orientation, la mémoire et l’incidence de l’état enrichi ; ses réalisations de spin et de jauge se composent sur les mêmes fibres matérielles. La composante électromagnétique appartient à la réalisation électrofaible : elle n’est pas introduite comme une interaction supplémentaire qui en serait déconnectée. L’action totale utilise les coefficients et les lecteurs déjà construits, dont la section de Planck, la longueur radiale et la fonctionnelle de Barbero--Immirzi.

Le contenu de cette intégration ne se limite pas à réunir des termes dans une expression. La variation de la même action produit le courant total de spin et la source métrique conjointe. L’élimination de la contorsion conserve les termes croisés entre espèces ; la transformation de Legendre conserve les moments matériels et produit les contraintes totales. Les preuves des sections~\ref{hmtcuatro:gea:seccion} et~\ref{hmtcuatro:accion-retroaccion-restricciones} établissent la rétroaction, le caractère de première classe et la propagation de ces contraintes dans la réalisation régulière spécifiée.

En particulier, soit $\Theta_{\rm tot}$ le potentiel canonique total obtenu dans cette transformation, $\mathcal Z=\{C_A=0\}$ la surface régulière des contraintes et $X_N=N^AX_{C_A}$, $X_M=M^BX_{C_B}$ deux déformations caractéristiques admissibles. La connexion induite par l’action possède $H_N^{\rm tot,can}=-\Theta_{\rm tot}(X_N)$. Sa courbure, calculée sans omettre les dérivées des coefficients de déformation, est
\begin{equation}
\begin{aligned}
\mathcal F^{\rm tot,can}_{N,M}
&=\delta_NH_M^{\rm tot,can}-\delta_MH_N^{\rm tot,can}
 -H_{[X_N,X_M]}^{\rm tot,can}
 +\frac{i}{\hbar}[H_N^{\rm tot,can},H_M^{\rm tot,can}]\\
&=-\mathrm d\Theta_{\rm tot}(X_N,X_M)
 =N^AM^B f_{AB}{}^{D}C_D.
\end{aligned}
\label{hmtcuatro:conclusion-curvatura-total}
\end{equation}
Par conséquent,
\[
 \left.\mathcal F^{\rm tot,can}_{N,M}\right|_{\mathcal Z}=0.
\]
La formule inclut la déformation tangentielle et les actions de jauge du crochet total ; si l’on travaille dans le quotient, elle exprime la même égalité modulo ces directions verticales. La représentation préquantique conserve cette clôture au moyen de l’identité de commutateurs démontrée dans \eqref{hmtcuatro:representacion-precuantica}.

Ainsi, l’incorporation gravitationnelle fait partie de la même action, de ses sources et de son algèbre de contraintes ; ce n’est pas une équation juxtaposée après la construction des autres interactions. L’annulation démontrée est celle de la courbure caractéristique canonique ; elle n’exige pas l’annulation de la courbure de l’espace-temps ni des intensités des champs de jauge. Son domaine conserve le corepère non dégénéré, la carte régulière et les conditions de bord utilisées dans la preuve. Elle ne supprime pas non plus la mémoire ni la monodromie globale des histoires qui ont donné lieu à la réalisation.
